% PAGES: 31-38, 39-45

\section{Nonsingular matrices}

From a numerical perspective, a nonsingular matrix $A$ is a square matrix
such that any linear system of the form $Ax=b$ has a unique solution $x$.
Solving linear systems corresponding to nonsingular matrices is often
required in financial applications and is a fundamental numerical linear
algebra problem.

The formal definition of a nonsingular matrix is given below:

\begin{definition}
The square matrix $A$ is nonsingular if and only if there exists another
matrix, called the inverse matrix of $A$ and denoted by $A^{-1}$, such
that

\begin{equation}
AA^{-1} \;=\; A^{-1}A \;=\; I.
\end{equation}
\end{definition}

\begin{definition}
The square matrix $A$ is singular if and only if the matrix $A$ is not
nonsingular.
\end{definition}

From (1.73), we obtain that the matrix $A$ satisfies the definition of the
inverse of the matrix $A^{-1}$. Thus, the inverse of the inverse of a
nonsingular matrix is equal to the matrix itself, i.e.,

\begin{equation}
\left(A^{-1}\right)^{-1} \;=\; A.
\end{equation}

Note that, for Definition 1.12 to make sense, the inverse matrix of a
nonsingular matrix as given by (1.73) must be unique, which is indeed the
case: Assume that $B_1$ and $B_2$ are inverse matrices of a nonsingular
matrix $A$. We will show that $B_1=B_2$. From (1.73), we find that

\[
AB_1 \;=\; B_1A \;=\; I \quad\text{and}\quad AB_2 \;=\; B_2A \;=\; I.
\]

By multiplying $I=AB_1$ to the left by $B_2$ and using the fact that
$B_2A=I$, we obtain that

\[
B_2 \;=\; B_2(AB_1) \;=\; (B_2A)B_1 \;=\; IB_1 \;=\; B_1.
\]

\begin{theorem}
Let $A$ be a square matrix of size $n$. The conditions below can be
regarded as equivalent definitions for $A$ to be a nonsingular matrix:

\begin{align}
A\ \text{nonsingular} &\iff \mathrm{Null}(A)=\{0\} \\
&\iff \mathrm{Range}(A)=\R^n \\
&\iff \rank(A)=n.
\end{align}

Similarly, the conditions below\footnote{Other necessary and sufficient
conditions for the matrix $A$ to be singular are
\[
A\ \text{singular} \iff \mathrm{Range}(A)\ne\R^n \iff \rank(A)<n.
\]
These conditions have less practical value, and are not included in
Theorem 1.1.} can be regarded as equivalent definitions for $A$ to be a
singular matrix:

\begin{align}
A\ \text{singular} &\iff \mathrm{Null}(A)\ne\{0\} \\
&\iff \text{there exists}\ v\in\R^n,\ v\ne 0,\ \text{such that}\ Av=0.
\end{align}
\end{theorem}

The proof of Theorem 1.1 is of less relevance for our purposes and can be
found, e.g., in Strang \cite{41}.

A very useful equivalent characterization of a matrix being nonsingular
can be given in terms of the determinant of the matrix; see section
10.1.1 for details.

\begin{theorem}
Let $A$ be a square matrix. Then,

\begin{align}
A\ \text{nonsingular} &\iff \det(A)\ne 0; \\
A\ \text{singular} &\iff \det(A)=0.
\end{align}
\end{theorem}

\begin{lemma}
Let $A$ and $B$ be square matrices of the same size. If $AB=I$, then $A$
and $B$ are nonsingular matrices, $B$ is the inverse matrix of $A$ (and
$A$ is the inverse matrix of $B$), and therefore $BA=I$.
\end{lemma}

\begin{proof}
Recall that the determinant of the identity matrix is 1, i.e.,
$\det(I)=1$; cf.\ (10.1). Since $AB=I$, it follows from Lemma 10.1 that
$\det(A)\det(B)=1$, and therefore both $\det(A)\ne 0$ and $\det(B)\ne 0$.
From (1.80), we conclude that $A$ and $B$ are nonsingular matrices. Let
$A^{-1}$ be the inverse of $A$. Multiplying $AB=I$ to the left by
$A^{-1}$ and using the fact that $A^{-1}A=I$, we find that

\[
A^{-1} \;=\; A^{-1}(AB) \;=\; (A^{-1}A)B \;=\; B.
\]

We conclude that $B$ is the inverse of $A$, and therefore
$BA=A^{-1}A=I$.
\end{proof}

\begin{lemma}
Let $A$ and $B$ be square matrices of the same size.

(i) The matrix $AB$ is nonsingular if and only if both $A$ and $B$ are
nonsingular.

(ii) If $A$ and $B$ are nonsingular, then\footnote{Note that (1.82)
extends as follows: $\left(\prod_{i=1}^{p}A_i\right)^{-1} =
\prod_{i=1}^{p}A_{p+1-i}^{-1}$. A proof can be given by induction; see an
exercise at the end of the chapter.}

\begin{equation}
(AB)^{-1} \;=\; B^{-1}A^{-1}.
\end{equation}
\end{lemma}

\begin{proof}
(i) From (1.80), we find that $AB$ is nonsingular if and only if
$\det(AB)\ne 0$, i.e., if and only if $\det(A)\det(B)\ne 0$; cf.\ (10.9).
Thus, the matrix $AB$ is nonsingular if and only if $\det(A)\ne 0$ and
$\det(B)\ne 0$. Using again (1.80), we conclude that the matrix $AB$ is
nonsingular if and only if both matrices $A$ and $B$ are nonsingular.

(ii) If $A$ and $B$ are nonsingular, it follows that

\[
\left(B^{-1}A^{-1}\right)AB \;=\; B^{-1}\left(A^{-1}A\right)B \;=\; B^{-1}B \;=\; I,
\]

since $A^{-1}A=I$ and $B^{-1}B=I$. From Definition 1.12, we conclude that
$(AB)^{-1}=B^{-1}A^{-1}$.
\end{proof}

It is important to note that the inverse of the product of two matrices
is \textit{not} the product of the inverses, i.e.,

\[
(AB)^{-1} \;\ne\; A^{-1}B^{-1}.
\]

Also, note the similarity between (1.82), i.e., $(AB)^{-1}=B^{-1}A^{-1}$,
and the result $(AB)^t=B^tA^t$ for the transposition of the product of
two matrices; cf.\ (1.24).

\begin{lemma}
Let $A$ be a square matrix. If $A$ is nonsingular, then $A^t$ is a
nonsingular matrix and

\begin{equation}
\left(A^t\right)^{-1} \;=\; \left(A^{-1}\right)^t.
\end{equation}
\end{lemma}

\begin{proof}
Note that it is enough to show that $\left(A^{-1}\right)^tA^t=I$; cf.\
Lemma 1.6. This follows from (1.24), since

\[
\left(A^{-1}\right)^tA^t \;=\; \left(AA^{-1}\right)^t \;=\; I^t \;=\; I. \qedhere
\]
\end{proof}

\begin{lemma}
The inverse of a symmetric nonsingular matrix is a symmetric matrix.

In other words, if $A$ is a symmetric nonsingular matrix, then
\[
\left(A^{-1}\right)^{t} \;=\; A^{-1}.
\]
\end{lemma}

\begin{proof}
If $A$ is a symmetric matrix, then $A^{t}=A$. From (1.83), we find that
\[
\left(A^{-1}\right)^{t} \;=\; \left(A^{t}\right)^{-1} \;=\; A^{-1}. \qedhere
\]
\end{proof}
